% GENERATED FILE. Edit canonical lessons, then run npm run generate:books.
% canonical-source-hash: fnv1a64:bd1d6083429ce555
% canonical-lessons: PT-C69-verde, PT-C69-cinzento, PT-C69-doce, PT-C69-limpo, PT-C69-sujo

\chapter{Green, Grey, Sweet, Clean, Dirty}
\label{ch:pt-green-grey-sweet-clean-dirty}
\hlchaptermodality{69}

\begin{chapteropening}
\textbf{By the end of this chapter:} \emph{I can say green, grey, sweet, clean and dirty.}

Complete the last lesson of chapter 69: I can say green, grey, sweet, clean and dirty.
\end{chapteropening}

\section[verde]{verde — green}

\label{lesson:PT-C69-verde}



\begin{quote}
Before the new one: say the Portuguese for poor, then the Portuguese for yellow.
\end{quote}

\subsection*{You'll want to know: verde}
\textbf{verde} — \textquotedblleft{}green\textquotedblright{}.

\begin{cousinweb}
\textbf{verde} $\leftarrow$ Latin \emph{viridis}, \textquotedblleft{}green\textquotedblright{} (EtymDB 2.1), shortened in speech to \emph{virdis} — inherited.
\end{cousinweb}

\begin{grammarlens}[title={What you've built}]
One more word for this chapter.
\end{grammarlens}

\subsection*{Your turn}
\begin{itemize}
\raggedright
  \item \emph{Say it:} \emph{verde}
  \item \emph{Say it:} \emph{verde}, once more
  \item \emph{Say it:} say \emph{verde}
  \item \emph{Recall:} say the Portuguese for poor, then the Portuguese for yellow, then say \emph{verde} again
\end{itemize}

\begin{tcolorbox}[breakable,colback=teal!4,colframe=teal!35!black,title={Before you move on}]
What does verde mean? (\textquotedblleft{}Green\textquotedblright{}.) Say it once more.
\end{tcolorbox}
\section[cinzento]{cinzento — grey}

\label{lesson:PT-C69-cinzento}



\begin{quote}
Before the new one: say the Portuguese for yellow, then the Portuguese for green.
\end{quote}

\subsection*{You'll want to know: cinzento}
\textbf{cinzento} — \textquotedblleft{}grey\textquotedblright{}.

\begin{cousinweb}
\textbf{cinzento} is \textquotedblleft{}ash-coloured\textquotedblright{}: built from the word for ash and \emph{-ento} (EtymDB 2.1). The ash word $\leftarrow$ Vulgar Latin \emph{*cinisia}, from Latin \emph{cinis}, \textquotedblleft{}ash\textquotedblright{} (EtymDB 2.1).
\end{cousinweb}

\begin{grammarlens}[title={What you've built}]
One more word for this chapter.
\end{grammarlens}

\subsection*{Your turn}
\begin{itemize}
\raggedright
  \item \emph{Say it:} \emph{cinzento}
  \item \emph{Say it:} \emph{cinzento}, once more
  \item \emph{Say it:} say \emph{cinzento}
  \item \emph{Recall:} say the Portuguese for yellow, then the Portuguese for green, then say \emph{cinzento} again
\end{itemize}

\begin{tcolorbox}[breakable,colback=teal!4,colframe=teal!35!black,title={Before you move on}]
What does cinzento mean? (\textquotedblleft{}Grey\textquotedblright{}.) Say it once more.
\end{tcolorbox}
\section[doce]{doce — sweet}

\label{lesson:PT-C69-doce}



\begin{quote}
Before the new one: say the Portuguese for green, then the Portuguese for grey.
\end{quote}

\subsection*{You'll want to know: doce}
\textbf{doce} — \textquotedblleft{}sweet\textquotedblright{}.

\begin{cousinweb}
\textbf{doce} $\leftarrow$ Latin \emph{dulcis}, \textquotedblleft{}sweet\textquotedblright{} (EtymDB 2.1) — inherited.
\end{cousinweb}

\begin{grammarlens}[title={What you've built}]
One more word for this chapter.
\end{grammarlens}

\subsection*{Your turn}
\begin{itemize}
\raggedright
  \item \emph{Say it:} \emph{doce}
  \item \emph{Say it:} \emph{doce}, once more
  \item \emph{Say it:} say \emph{doce}
  \item \emph{Recall:} say the Portuguese for green, then the Portuguese for grey, then say \emph{doce} again
\end{itemize}

\begin{tcolorbox}[breakable,colback=teal!4,colframe=teal!35!black,title={Before you move on}]
What does doce mean? (\textquotedblleft{}Sweet\textquotedblright{}.) Say it once more.
\end{tcolorbox}
\section[limpo]{limpo — clean}

\label{lesson:PT-C69-limpo}



\begin{quote}
Before the new one: say the Portuguese for grey, then the Portuguese for sweet.
\end{quote}

\subsection*{You'll want to know: limpo}
\textbf{limpo} — \textquotedblleft{}clean\textquotedblright{}.

\begin{cousinweb}
\textbf{limpo} $\leftarrow$ Latin \emph{limpidus}, \textquotedblleft{}clear\textquotedblright{} (EtymDB 2.1) — inherited. It lost a whole syllable, as inherited words do.
\end{cousinweb}

\begin{grammarlens}[title={What you've built}]
One more word for this chapter.
\end{grammarlens}

\subsection*{Your turn}
\begin{itemize}
\raggedright
  \item \emph{Say it:} \emph{limpo}
  \item \emph{Say it:} \emph{limpo}, once more
  \item \emph{Say it:} say \emph{limpo}
  \item \emph{Recall:} say the Portuguese for grey, then the Portuguese for sweet, then say \emph{limpo} again
\end{itemize}

\begin{tcolorbox}[breakable,colback=teal!4,colframe=teal!35!black,title={Before you move on}]
What does limpo mean? (\textquotedblleft{}Clean\textquotedblright{}.) Say it once more.
\end{tcolorbox}
\section[sujo]{sujo — dirty}

\label{lesson:PT-C69-sujo}



\begin{quote}
Before the new one: say the Portuguese for sweet, then the Portuguese for clean.
\end{quote}

\subsection*{You'll want to know: sujo}
\textbf{sujo} — \textquotedblleft{}dirty\textquotedblright{}.

\begin{cousinweb}
\textbf{sujo} $\leftarrow$ Latin \emph{sucidus}, \textquotedblleft{}juicy, greasy\textquotedblright{} (EtymDB 2.1), from \emph{sucus}, \textquotedblleft{}juice\textquotedblright{} — inherited.
\end{cousinweb}

\begin{grammarlens}[title={What you've built}]
One more word for this chapter.
\end{grammarlens}

\subsection*{Your turn}
\begin{itemize}
\raggedright
  \item \emph{Say it:} \emph{sujo}
  \item \emph{Say it:} \emph{sujo}, once more
  \item \emph{Say it:} say \emph{sujo}
  \item \emph{Recall:} say the Portuguese for sweet, then the Portuguese for clean, then say \emph{sujo} again
\end{itemize}

\begin{tcolorbox}[breakable,colback=teal!4,colframe=teal!35!black,title={Before you move on}]
What does sujo mean? (\textquotedblleft{}Dirty\textquotedblright{}.) Say it once more.
\end{tcolorbox}
